\section{Methodology}
\label{sec:method}

We studied two offerings of a project-based undergraduate software-engineering
course at a large public university in South America: 2025.2 ($N=9$
analyzable teams) and 2026.1 ($N=5$), for $N=14$ team-semesters and
approximately 70 participating students. Teams of roughly five students had
autonomy over process model and project management approach. Course records
did not formally record teams' process selections or adherence to ceremonies,
so we cannot report a distribution across Scrum, Kanban, other Agile variants,
or unguided work. Each team was required to build an operational software
project with an integrated GenAI component, but no AI architecture was
mandated. Retrieval-Augmented Generation (RAG) was presented as one possible
architecture~\cite{lewis2020retrieval}. The analyzed systems spanned
healthcare, finance, education, research, ecotourism, onboarding, and tutoring
domains, and used project-specific integrations such as summarization,
document/profile analysis, chatbots, recommendations, and content generation.

We report this observational, multi-method study with reference to established
empirical software-engineering research and case-study guidance
\cite{stol2018abc,runeson2009case}.

\label{sec:data-sources}
The study uses nominal $T_1$--$T_3$ markers, but survey and evaluator dates do
not coincide. Survey forms place $T_1$ on September 5, 2025, and March 5--6,
2026. Taking the earliest $T_1$ response date as course day one solely to
estimate elapsed weeks, the survey waves fall approximately in weeks 1, 9, and
13 for 2025.2, and weeks 1, 11, and 15 for 2026.1. Evaluator scorecards were
collected after status presentations on October 17 and 24, November 14 and 21,
and December 5 and 12 in 2025.2, and on April 24, May 22, and June 19 in
2026.1. Relative to the same anchor, these dates fall approximately in weeks
7--8, 11--12, and 14--15 for 2025.2, and weeks 8, 12, and 16 for 2026.1.
These are date-derived estimates, not official course-week labels. The
versioned records do not define a shared artifact checklist for each
checkpoint or document a repository-onboarding change that explains the
increase in observed $T_1$ architectural artifacts from 4/9 teams in 2025.2 to
5/5 in 2026.1; we therefore report that difference descriptively without
attributing a cause.

The study combines four evidence streams. First, complete Git histories
provide commits, line changes, file provenance, branch information, and
repository-visible artifacts. Students were instructed to keep requirements,
architecture, documentation, and source code in the repository, but Git
captures only repository-visible evidence--artifacts and changes committed to
version control--rather than AI-use telemetry such as prompt logs, API calls,
or IDE acceptance events. Second, independent evaluator scorecards were
collected at $T_1$--$T_3$ through blind scoring after status presentations,
with 4 evaluators in 2025.2 and 3 in 2026.1. Third, multi-checkpoint surveys
captured demographics, experience, GenAI adoption, perceived role disruption,
and project expectations at cohort/checkpoint grain; survey responses carry no
team identifier. Fourth, a bounded 2025.2 transcript corpus covers status
presentations and supervisor discussions, segmented into 100 chunks across
$T_1$, $T_2$, and $T_3$. Transcript evidence is a global checkpoint corpus
and is not imputed for 2026.1.

All raw streams were anonymized before analysis. Analytical artifacts exclude
direct personal identifiers; restricted provenance records preserve prompts,
responses, hashes, and usage metadata for auditability.
% PROVENANCE: Canonical inputs and checksums are listed in
% paper_v9/data/manifests/input_inventory.json; survey and evaluator schemas
% are validated by paper_v9/tests/test_input_contracts.py.

\subsection{Analytical Approach}

The analysis is descriptive and preserves the scale and grain of each evidence
stream. Means are reported with standard deviations for approximately
interval-like scores; medians and interquartile ranges are used for skewed
counts and ratios. Proportions retain denominators. Rolling seven-day windows
are indexed relative to evaluator anchors and are temporal summaries, not
independent observations. Exploratory associations use Spearman's rank
correlation $\rho$~\cite{spearman1904proof}, which tests whether two variables
tend to move monotonically together after replacing raw values by ranks. This
choice avoids assuming linear relationships or normally distributed variables
in a small, mixed-scale panel. Leave-one-out analyses recompute associations
after removing each eligible team-semester.

In addition to the primary M1--M9 families, robustness views include
non-overlapping phase bins, final-seven-day concentration contrasts by score
change and planning scope, operational-regularity proxies,
technical-complexity profiles, and leave-one-out sensitivity ranges.
\textit{Epistemological boundaries and grain separation:} Because this study
observes $N=14$ team-semesters and approximately 70 students without a non-AI
control group or IDE-level telemetry, three principles govern interpretation
of the Section IV results. First, \textit{strict grain preservation}: student-
level survey waves (M1/M2), global checkpoint transcripts (M5), and
team-semester repository/scorecard metrics (M3, M4, M6--M9) retain distinct
denominators (Table~\ref{tab:metrics_summary}) and are never pooled or
cross-correlated across incompatible grains. Second, \textit{exploratory
inference}: Spearman correlations ($\rho$), $p$-values, phase bins,
final-seven-day contrasts, operational-regularity proxies,
technical-complexity stratifications, and leave-one-out ranges (which show
sensitivity to influential cases rather than confidence intervals) are
descriptive or robustness diagnostics, not confirmatory tests or causal
estimators. Third,
\textit{proxy transparency}: missing repository or transcript observations
remain explicit observation states rather than being zero-filled or imputed as
low quality. These boundaries are stated here and in Section~VI rather than
repeated for every metric.

\subsection{Metric Families}

RQ1 is operationalized at cohort/checkpoint grain. Self-Reported AI Usage and
Perceptions (M1) separates six perception families from self-reported usage
frequency, tasks, tools, prior experience, autonomy/tool balance, and
exploratory career topics. Perceived Role Disruption (M2) preserves
role-specific Likert distributions and paired sensitivity summaries. Neither
metric is behavioral AI-use telemetry.

RQ2 uses repository and transcript evidence with explicit grain separation.
Repository Activity (M3) summarizes commit timing, author participation,
rolling windows, final-seven-day concentration, and non-overlapping phase
shares for 14 team-semesters. Clean-Change Dynamics (M4) applies a versioned
clean-path policy and reports clean churn separately from total churn.
Coordination-Language Evidence (M5) reconstructs one 2025.2 global transcript
corpus per checkpoint and deterministically counts versioned lexical markers;
it is not correlated with team-semester repository observations.

RQ3 uses team-semester artifacts and outcomes. Upfront Architecture Evidence
(M6) combines repository-visible T1 artifact presence/scope (M6a) with
structured Planning Content Evidence (M6b) from T1 commit subjects for nine
observed team-semesters; five cases without T1 subjects remain unavailable.
Repository Inactivity (M7) reports timing and coverage diagnostics, not
planning omission. Clean-Path Rework (M8) reports late changes to previously
observed clean source/test paths, using file provenance as a rework proxy while
preserving baseline eligibility. Artifact--Rework--
Outcome Associations (M9) reports exploratory Spearman associations among T1
artifact predictors, M8 rework, and independent T3 scorecard outcomes, with
ratio analyses restricted to baseline-eligible cases. Technical-complexity
profiles contextualize rework and concentration, but are not causal
adjustments.

\begin{table}[t]
\caption{Summary of Metric Families, Operational Grains, and Coverage}
\label{tab:metrics_summary}
\centering
\small
\begin{tabular}{@{}lp{4.1cm}p{3.0cm}@{}}
\toprule
\textbf{ID} & \textbf{Metric Family and Components} & \textbf{Primary Grain ($N$)} \\
\midrule
M1 & Self-Reported AI Usage \& Perceptions (6 perception scales + task/tool frequency) & Cohort $\times$ checkpoint ($6$ cells) \\
\midrule
M2 & Perceived Role Disruption (5-point Likert across 6 SE roles + paired shifts) & Cohort $\times$ checkpoint $\times$ role \\
\midrule
M3 & Repository Activity (commit timing, rolling windows, author Gini, phase bins) & Team-semester / window ($N=14$) \\
\midrule
M4 & Clean-Change Dynamics (source/test line churn under clean-path policy) & Team-semester / checkpoint ($N=14$) \\
\midrule
M5 & Coordination-Language Evidence (deterministic lexical markers per 1k tokens) & Global corpus / checkpoint (\texttt{2025.2}) \\
\midrule
M6 & Upfront Architecture Evidence: artifact scope (M6a, $N=14$) \& commit-subject categories (M6b, $n=9$) & Team-semester ($N=14$ / $n=9$) \\
\midrule
M7 & Repository Inactivity (7-day window coverage diagnostic) & Team-semester / window ($N=14$) \\
\midrule
M8 & Clean-Path Rework (late modifications to prior clean source/test files) & Team-semester ($N=14$; $n=12$ ratio) \\
\midrule
M9 & Artifact--Rework--Outcome Associations (Spearman $\rho$ + leave-one-out sensitivity) & Team-semester / stratum ($N=14$) \\
\bottomrule
\end{tabular}
\end{table}

% PROVENANCE: Family/component definitions and exact artifact paths are indexed
% in paper_v9/data/results/results_summary.json. Metric-level transformations
% are documented in paper_v9/metricas_v8-analysis/analyze_m1.md through
% analyze_m9.md and in the corresponding verification notebooks.

\subsection{Language-Model Assistance}

Language-model assistance is used for bounded data extraction and is kept
separate from deterministic metric construction. M1 uses coded survey
responses and M5 uses deterministic lexical coding; neither receives a new
language-model score. M6b uses structured extraction from $T_1$ commit subjects
to identify goals, architecture/design, task decomposition, risk/dependency,
and literal evidence quotes. Outputs were reviewed by the research team and
retained as separate exploratory categories, never as a composite
planning-quality or architectural-quality score.
